% GENERATED FILE. Edit canonical lessons, then run npm run generate:books.
% canonical-source-hash: fnv1a64:401c36870e025e79
% canonical-lessons: MW-C155-jau-javai, MW-C155-khau-pivai, MW-C155-dudh, MW-W46-dha, MW-R155-eat-drink-go-come

\chapter{Eat, Drink, Go, Come}
\label{ch:mw-a2-eat-drink-go-come}
\hlchaptermodality{155}

\begin{chapteropening}
\textbf{By the end of this chapter:} \emph{I can say what I and others eat, drink, and where we go and come from, every day, and write the letter \mw{ध} from \mw{दूध}.}

Complete MW-R155-eat-drink-go-come: use \mw{खा}, \mw{पी}, \mw{जा} and \mw{आ} in the present, name milk, and write \mw{ध}.
\end{chapteropening}

\section[\texorpdfstring{\mw{जाऊं} · \mw{जावै} (jāū̃ · jāvai)}{जाऊं · जावै (jāū̃ · jāvai)}]{\mw{हूं} \mw{जाऊं} \mw{हूं} — a stem that ends in a vowel}

\label{lesson:MW-C155-jau-javai}



\begin{quote}
Say \emph{to go}, \textbf{\mw{जाणो}}. Then \emph{I listen to songs} and \emph{he listens to songs}. And \emph{the market}, \textbf{\mw{बाजार}}.
\end{quote}

\begin{grammarlens}[title={Grammar Lens: \mw{जा}}]
\textbf{\mw{जाणो}} leaves \textbf{\mw{जा}}, a stem that ends in a vowel. For \emph{I}, its ending is the full letter \textbf{\mw{ऊं}}:

\begin{quote}
\textbf{\mw{हूं} \mw{रोज} \mw{बाजार} \mw{जाऊं} \mw{हूं।}}
\end{quote}

\begin{quote}
\emph{hū̃ roj bājār jāū̃ hū̃.} — I go to the market every day.
\end{quote}

For \textbf{\mw{वो}}, a \textbf{\mw{व}} comes in before \textbf{-\mw{ै}}:

\begin{quote}
\textbf{\mw{वो} \mw{रोज} \mw{बाजार} \mw{जावै} \mw{है।}}
\end{quote}

\begin{quote}
\emph{vo roj bājār jāvai hai.} — He goes to the market every day.
\end{quote}

After \textbf{\mw{आप}}, the \textbf{\mw{व}} stays: \textbf{\mw{आप} \mw{जावो} \mw{हो}}, the ending of the request \textbf{\mw{लावो}}. Old Marwari had \emph{jāvaï}, \textquotedblleft{}goes\textquotedblright{} (CDIAL 10452).
\end{grammarlens}

\subsection*{Your turn}
Say these aloud:

\begin{itemize}
\raggedright
  \item \textbf{\mw{हूं} \mw{रोज} \mw{बाजार} \mw{जाऊं} \mw{हूं।}} — \emph{hū̃ roj bājār jāū̃ hū̃.} — I go to the market every day.
  \item \textbf{\mw{मां} \mw{रोज} \mw{बाजार} \mw{जावै} \mw{है।}} — \emph{mā̃ roj bājār jāvai hai.} — Mother goes to the market every day.
  \item \textbf{\mw{आप} \mw{बाजार} \mw{जावो} \mw{हो।}} — \emph{āp bājār jāvo ho.} — You go to the market.
\end{itemize}

\begin{tcolorbox}[breakable,colback=teal!4,colframe=teal!35!black,title={Before you move on}]
\emph{I go}? (\textbf{\mw{जाऊं} \mw{हूं}}.) \emph{He goes}? (\textbf{\mw{जावै} \mw{है}}.) \emph{You go}, to \textbf{\mw{आप}}? (\textbf{\mw{जावो} \mw{हो}}.)
\end{tcolorbox}
\section[\texorpdfstring{\mw{खाऊं} · \mw{पीवै} · \mw{आवै} (khāū̃ · pīvai · āvai)}{खाऊं · पीवै · आवै (khāū̃ · pīvai · āvai)}]{\mw{खा}, \mw{पी}, \mw{आ} — three more vowel stems}

\label{lesson:MW-C155-khau-pivai}



\begin{quote}
Say \emph{I go to the market every day}. Then \emph{to eat}, \emph{to drink}, \emph{to come}: \textbf{\mw{खाणो}}, \textbf{\mw{पीणो}}, \textbf{\mw{आवणो}}.
\end{quote}

\begin{grammarlens}[title={Grammar Lens: the same three endings}]
\textbf{\mw{खा}}, \textbf{\mw{पी}} and \textbf{\mw{आ}} end in vowels, and go exactly like \textbf{\mw{जा}}:

\begin{quote}
\textbf{\mw{हूं} \mw{रोज} \mw{दिनूगै} \mw{चाय} \mw{पीऊं} \mw{हूं।}}
\end{quote}

\begin{quote}
\emph{hū̃ roj dinūgai chāy pīū̃ hū̃.} — I drink tea early every morning.
\end{quote}

\begin{quote}
\textbf{\mw{राम} \mw{रोटी} \mw{खावै} \mw{है।}}
\end{quote}

\begin{quote}
\emph{rām roṭī khāvai hai.} — Ram eats roti.
\end{quote}

\begin{quote}
\textbf{\mw{आप} \mw{रोज} \mw{आवो} \mw{हो।}}
\end{quote}

\begin{quote}
\emph{āp roj āvo ho.} — You come every day.
\end{quote}

\textbf{\mw{आवणो}} leaves \textbf{\mw{आव}}, but its present follows \textbf{\mw{जा}}: \textbf{\mw{आऊं}}, \textbf{\mw{आवै}}, \textbf{\mw{आवो}}. Old Marwari had \emph{āvaï}, \textquotedblleft{}comes\textquotedblright{} (CDIAL 1200).
\end{grammarlens}

\subsection*{Your turn}
Every cell of the present, now that each one is taught:

\noindent
\begin{tabularx}{\linewidth}{@{}>{\raggedright\arraybackslash}X>{\raggedright\arraybackslash}X>{\raggedright\arraybackslash}X@{}}
\toprule
\textbf{who} & \textbf{\mw{बोल}-} & \textbf{\mw{जा}-} \\
\midrule
\textbf{\mw{हूं}} & \textbf{\mw{बोलूं} \mw{हूं}} & \textbf{\mw{जाऊं} \mw{हूं}} \\
\textbf{\mw{वो}} & \textbf{\mw{बोलै} \mw{है}} & \textbf{\mw{जावै} \mw{है}} \\
\textbf{\mw{आप}} & \textbf{\mw{बोलो} \mw{हो}} & \textbf{\mw{जावो} \mw{हो}} \\
\bottomrule
\end{tabularx}

Say these aloud:

\begin{itemize}
\raggedright
  \item \textbf{\mw{हूं} \mw{रोटी} \mw{खाऊं} \mw{हूं।}} — \emph{hū̃ roṭī khāū̃ hū̃.} — I eat roti.
  \item \textbf{\mw{वो} \mw{चाय} \mw{पीवै} \mw{है।}} — \emph{vo chāy pīvai hai.} — He drinks tea.
  \item \textbf{\mw{वो} \mw{रोज} \mw{आवै} \mw{है।}} — \emph{vo roj āvai hai.} — He comes every day.
\end{itemize}

\begin{tcolorbox}[breakable,colback=teal!4,colframe=teal!35!black,title={Before you move on}]
\emph{He eats}? (\textbf{\mw{खावै} \mw{है}}.) \emph{I come}? (\textbf{\mw{आऊं} \mw{हूं}}.) \emph{You drink}, to \textbf{\mw{आप}}? (\textbf{\mw{पीवो} \mw{हो}}.)
\end{tcolorbox}
\section[\texorpdfstring{\mw{दूध} (dūdh)}{दूध (dūdh)}]{\mw{दूध} (dūdh) — milk}

\label{lesson:MW-C155-dudh}



\begin{quote}
Before the new one: say \emph{tea}, \emph{water} and \emph{yoghurt}. (\textbf{\mw{चाय}}, \textbf{\mw{पाणी}}, \textbf{\mw{दही}}.)
\end{quote}

\subsection*{You'll want to know: \mw{दूध}}
\textbf{\mw{दूध}} — \emph{dūdh} — \textquotedblleft{}milk\textquotedblright{}.

\begin{quote}
\textbf{\mw{हूं} \mw{रोज} \mw{दूध} \mw{पीऊं} \mw{हूं।}}
\end{quote}

\begin{quote}
\emph{hū̃ roj dūdh pīū̃ hū̃.} — I drink milk every day.
\end{quote}

It ends on a letter you have not written yet: \emph{dh}, a \textbf{\mw{द}} with a breath after it. The next lesson puts it in your hand.

\begin{cousinweb}
\textbf{\mw{दूध}} (\emph{dūdh}) is inherited from Sanskrit \emph{dugdha}, \textquotedblleft{}milk\textquotedblright{}, a word built on the verb for milking. CDIAL lists Marwari \emph{dūdh}, masculine (CDIAL 6391).
\end{cousinweb}

\begin{grammarlens}[title={What you've built}]
\textbf{\mw{दूध}} is masculine: \textbf{\mw{दूध} \mw{गरम} \mw{है।}} (\emph{dūdh garam hai.}), the milk is hot.
\end{grammarlens}

\subsection*{Your turn}
\begin{itemize}
\raggedright
  \item \emph{Say it:} \emph{dūdh}
  \item \emph{Say it:} \textbf{\mw{वो} \mw{दूध} \mw{पीवै} \mw{है।}} — \emph{vo dūdh pīvai hai.} — He drinks milk.
  \item \emph{Recall:} say \emph{tea}, then \emph{dūdh} again
\end{itemize}

\begin{tcolorbox}[breakable,colback=teal!4,colframe=teal!35!black,title={Before you move on}]
What does \textbf{\mw{दूध}} mean? (\textquotedblleft{}Milk\textquotedblright{}.) Say it once more.
\end{tcolorbox}
\section[\texorpdfstring{\mw{ध} (dha)}{ध (dha)}]{\mw{ध} (dha) — the letter that ends \mw{दूध}}

\label{lesson:MW-W46-dha}



\begin{quote}
Before the new one: write \textbf{\mw{द}} once, and name it.

Now say \textbf{\mw{दूध}} (\emph{dūdh}), \textquotedblleft{}milk\textquotedblright{}. It opens on \textbf{\mw{द}} and ends on a \textbf{\mw{द}} with a breath after it: \emph{dh}. This lesson puts that letter in your hand.
\end{quote}

\subsection*{Script you'll notice: \mw{ध}}
\textbf{\mw{ध}} — \emph{dha}, the breathed \emph{d}.

It ends \textbf{\mw{दूध}}. It sits beside \textbf{\mw{द}} as \textbf{\mw{घ}} sits beside \textbf{\mw{ग}}: the same place in the mouth, with a breath. On the page it is an upper spiral that ends in a shoulder, a lower bowl, and a right stem under the headline.

\subsection*{Writing: \mw{ध}}
\hlblockfigure{\detokenize{figures/MW-W46-dha-filmstrip.pdf}}{How \mw{ध} is written, stroke by stroke}

\begin{itemize}
\raggedright
  \item \textbf{1.} start at the upper spiral's inner crossing, curl around the small opening, widen left and down around the outer loop, then sweep right through the shoulder
  \item \textbf{2.} without lifting, turn back left along the shoulder to the left waist
  \item \textbf{3.} sweep down and around the lower bowl to its right junction
  \item \textbf{4.} climb up the right stem to the headline
  \item \textbf{5.} descend the right stem top-to-bottom
  \item \textbf{6.} lift and draw the top shirorekhā left-to-right
\end{itemize}

\textbf{Pen lifts: 1.}

\begin{quote}
The source shows four movements; native writers draw \mw{ध} as two pen-down strokes (81\% of HP Labs India's native-writer samples, 64 of 79). This ductus follows the four movements in order and direction, lifting only before the headline: after the shoulder the pen turns back along it into the lower bowl, then climbs the right stem before descending it.
\end{quote}

\begin{quote}
Stroke order is one attested teaching order; hands and teachers vary. Source: Opiaterein, ‘Deva-\mw{ध}-order.gif’, strokes 1–4, Wikimedia Commons, 10 May 2009.
\end{quote}

\subsection*{Your turn}
\begin{itemize}
\raggedright
  \item \emph{Look:} at \textbf{\mw{दूध}} and point at \mw{ध} at its end
  \item \emph{Trace it:} \mw{ध} three times, saying \emph{dh} as you finish each one
  \item \emph{Write it:} \mw{द}, then \mw{ध}, from memory
  \item \emph{Read it:} \textbf{\mw{दूध}} aloud once more
\end{itemize}

\begin{tcolorbox}[breakable,colback=teal!4,colframe=teal!35!black,title={Before you move on}]
Which letter is this — \mw{ध}? (\textbf{dha}, the breathed \emph{d}.) Which word did it come from? (\textbf{\mw{दूध}}, \emph{dūdh}, \textquotedblleft{}milk\textquotedblright{}.)
\end{tcolorbox}
\section[(dialogue)]{Eat, drink, go, come — a review}

\label{lesson:MW-R155-eat-drink-go-come}



\begin{quote}
Say \emph{I go to the market every day}.
\end{quote}

\subsection*{Your turn}
\begin{itemize}
\raggedright
  \item \emph{Say it:} \emph{I eat roti}; \emph{he eats roti}; \emph{you eat roti}, to \textbf{\mw{आप}}
  \item \emph{Say it:} \emph{I drink milk}; \emph{mother drinks milk every day}
  \item \emph{Say it:} \emph{I come every day}; \emph{he comes every day}
  \item \emph{Say it:} \emph{he goes to the market}; \emph{you go to the market}
  \item \emph{Read it:} \textbf{\mw{दूध}}, and name its last letter
  \item \emph{Read it:} \textbf{\mw{जाऊं}}, and name the letter after \textbf{\mw{जा}}
\end{itemize}

\begin{tcolorbox}[breakable,colback=teal!4,colframe=teal!35!black,title={Before you move on}]
\emph{I eat}? (\textbf{\mw{खाऊं} \mw{हूं}}.) \emph{He drinks}? (\textbf{\mw{पीवै} \mw{है}}.) \emph{You come}? (\textbf{\mw{आवो} \mw{हो}}.) \emph{Milk}? (\textbf{\mw{दूध}}.)
\end{tcolorbox}
